\documentclass[11pt, a4paper]{article}

\usepackage[utf8]{inputenc}
\usepackage[T1]{fontenc}
\usepackage{lmodern}
\usepackage{amsmath, amssymb, amsthm}
\usepackage[a4paper, margin=1in]{geometry}
\usepackage{microtype}
\usepackage[round]{natbib}
\usepackage{hyperref}
\usepackage{enumitem}
\usepackage{indentfirst}
\usepackage{xcolor}
\hypersetup{colorlinks=true, linkcolor=blue!50!black, urlcolor=blue!50!black, citecolor=blue!50!black}

\title{Wave Release Coordination under Vertical Resource Constraints in Multi-Story AMR Warehouses}
\author{Shiyue Hu \and (co-authors TBD)}
\date{May 2026}

\begin{document}

\maketitle

\section{Related Work}\label{sec:related}

We organize the relevant literature into four problem-class threads (warehouse OR sharing part of our setting) and three methodological threads (analytical tools bearing on our framework). Two recent surveys provide context: \citet{boysendekoster2025} retrospects fifty years of warehousing OR across three generations of distribution-center architecture, and \citet{boysenschneiderzulj2025} reviews fleet-coordination problems for electric vehicles in facility logistics, including AMR fleets.

\citet{lenoble2018} formulate order batching for vertical lift modules (VLMs), where each VLM is a single mechanical device and the decision is which item-orders to retrieve in one tray-load; the mathematics is internal to the device. Our setting differs at every scale: the vertical resource is a freight elevator carrying AMRs across building floors, the wave is a fleet-level co-occupation window, and multi-elevator capacity shared across a horizontal AMR fleet does not arise inside a single VLM.

\citet{chakravarty2025} compute provably optimal lift schedules for AMR fleets using Boolean satisfiability (SAT) in hotel and hospital deployments. Their decision space is downstream of ours: the SAT formulation takes robot--task assignments as exogenous and finds the time-optimal schedule. We instead treat the \emph{selection} of which orders enter the system together (wave composition) as an upstream tactical decision variable, with operational dispatch then taking that wave as input. Their SAT scaling caps near thirty agents per scenario; we target hundreds of orders per wave through a heuristic-friendly two-stage decomposition.

\citet{wu2024} address inbound job scheduling in four-way shuttle systems via a flexible flow-shop formulation, and \citet{keung2023} study order and pod assignment in multi-level robotic mobile fulfillment systems (RMFS). Both adopt a \emph{tier-captive} transport model: each shuttle is bound to a specific layer, with inter-layer transfers via dedicated lift--shuttle handover. Our setting has flexible AMRs queueing at shared elevators, producing a different bottleneck topology because AMRs from any floor compete concurrently for elevator service. This concurrency is what makes wave composition the binding tactical lever in our problem class but not in theirs.

The work most directly relevant to our tactical layer treats order grouping as a decision variable in planar (single-floor) settings without vertical resource coupling. \citet{scholz2017} formulate the \emph{Joint Order Batching, Assignment and Sequencing, and Picker Routing Problem} (JOBASRP) for multi-picker manual warehouses, showing the joint solution outperforms sequential decomposition by up to 84\%. \citet{zulj2022} extend this line to AMR-assisted picker-to-parts systems with a hybrid ALNS--tabu heuristic. \citet{qin2024} study multi-tote AMR fulfillment and notably \emph{use the term ``wave''} for processing batches, identifying an empirical optimum near 100 orders per wave in their single-layer setting; their ``wave'' is a \emph{cardinality} parameter, not a \emph{composition} decision on structured features. Our work extends this single-floor decision class to the multi-story setting where shared elevator capacity introduces a vertical-resource coupling absent from these planar precedents.

Beyond these problem-class neighbors, three methodological lines bear on our analytical apparatus, each addressing a precise gap one of our methodological contributions fills. Robust scheduling under model uncertainty \citep{wiesemann2013, lushen2021} handles parametric variation within a known model class, presupposing a single agreed-upon model family with an unknown parameter. Our elevator-modeling uncertainty is structurally different: throughput-aggregation and true-co-occupancy-batching are \emph{qualitatively distinct} models, not parameterizations of the same family. Our \textbf{Model-Dominance Hedge Rule} addresses this gap by collapsing structural model-class disagreement into a single closed-form decision under a per-wave dominance condition, with a worst-case loss bound that requires no online model selection.

The prediction-to-decision regret literature \citep{elmachtoub2022, vera2021} bounds algorithmic loss when a learned predictor drives a downstream optimizer. These bounds are \emph{training-time and pre-commitment}: they describe a specific algorithm's regret, not the \emph{post-hoc} gap between an oracle on a structured feature partition and a feature-informed policy on the same partition. Our \textbf{Bound-and-Gap framework} addresses this gap by formalizing that information gap, decomposable into a partition-intrinsic upper-tail term ($H_{\mathrm{up}}$) and a $\Phi$-policy slack term ($M_\Phi$).

Structural-feature decomposition of order properties has been treated via raw operational attributes: due dates and item locations \citep{scholz2017}, tour feasibility \citep{zulj2022}, and item popularity \citep{qin2024}. These characterizations are \emph{decision-driven} (capturing whether an order fits a tour) rather than \emph{structurally diagnostic} (decomposing why one wave is harder than another into orthogonal physical mechanisms), leaving operators without a way to localize the source of difficulty. Our \textbf{three-dimensional structured decomposition} $\Phi = (C, I, T)$, namely vertical concentration, directional imbalance, and temporal clustering, addresses this gap with physically interpretable per-unit coefficients along three orthogonal axes, each capturing a distinct vertical-resource stress.

None of the threads above treats the four problem-defining elements (wave composition as tactical decision variable, multi-story deployment, flexible AMR fleets, shared-elevator capacity coupling) as one coupled system: order-grouping work is planar; multi-story scheduling work is tier-captive; multi-agent robot--lift work has exogenous task selection; device-level VLM batching and the methodological literature each touch one piece. We formalize the joint problem class in Section~3 and develop our analytical tools ($\Phi$, the Bound-and-Gap framework, and the Model-Dominance Hedge Rule) in Section~4.

% Bibliography placeholder (use BibTeX in final draft)
\begin{thebibliography}{14}
  \bibitem[Boysen \& de Koster, 2025]{boysendekoster2025} Boysen, N., \& de Koster, R. (2025). 50 years of warehousing research: An operations research perspective. \emph{European Journal of Operational Research}, 320(3), 449--464. \url{https://doi.org/10.1016/j.ejor.2024.03.026}
  \bibitem[Boysen, Schneider \& \v{Z}ulj, 2025]{boysenschneiderzulj2025} Boysen, N., Schneider, M., \& \v{Z}ulj, I. (2025). Energy management for electric vehicles in facility logistics: A survey from an operational research perspective. \emph{European Journal of Operational Research} (Invited Review, in press). \url{https://doi.org/10.1016/j.ejor.2025.12.031}
  \bibitem[Chakravarty et al., 2025]{chakravarty2025} Chakravarty, A., Grey, M.\,X., Muthugala, M.\,A.\,V.\,J., \& Elara, R.\,M. (2025). Toward optimal multi-agent robot and lift schedules via Boolean satisfiability. \emph{Mathematics}, 13(18), 3031. \url{https://doi.org/10.3390/math13183031}
  \bibitem[Elmachtoub \& Grigas, 2022]{elmachtoub2022} Elmachtoub, A.\,N., \& Grigas, P. (2022). Smart ``Predict, then Optimize''. \emph{Management Science}, 68(1), 9--26. \url{https://doi.org/10.1287/mnsc.2020.3922}
  \bibitem[Keung et al., 2023]{keung2023} Keung, K.\,L., Lee, C.\,K.\,M., \& Ji, P. (2023). Assigning orders and pods to picking stations in a multi-level robotic mobile fulfillment system. \emph{Flexible Services and Manufacturing Journal}. \url{https://doi.org/10.1007/s10696-023-09491-0}
  \bibitem[Lenoble et al., 2018]{lenoble2018} Lenoble, N., Frein, Y., \& Hammami, R. (2018). Order batching in an automated warehouse with several vertical lift modules: Optimization and experiments with real data. \emph{European Journal of Operational Research}, 267(3), 958--976. \url{https://doi.org/10.1016/j.ejor.2017.12.037}
  \bibitem[Lu \& Shen, 2021]{lushen2021} Lu, M., \& Shen, Z.-J.\,M. (2021). A review of robust operations management under model uncertainty. \emph{Production and Operations Management}, 30(6), 1927--1943. \url{https://doi.org/10.1111/poms.13239}
  \bibitem[Qin et al., 2024]{qin2024} Qin, Z., Kang, Y., \& Yang, P. (2024). Making better order fulfillment in multi-tote storage and retrieval autonomous mobile robot systems. \emph{Transportation Research Part E: Logistics and Transportation Review}, 192, 103752. \url{https://doi.org/10.1016/j.tre.2024.103752}
  \bibitem[Scholz et al., 2017]{scholz2017} Scholz, A., Schubert, D., \& W\"ascher, G. (2017). Order picking with multiple pickers and due dates: Simultaneous solution of Order Batching, Batch Assignment and Sequencing, and Picker Routing Problems. \emph{European Journal of Operational Research}, 263(2), 461--478. \url{https://doi.org/10.1016/j.ejor.2017.04.038}
  \bibitem[Vera et al., 2021]{vera2021} Vera, A., Banerjee, S., \& Gurvich, I. (2021). Online allocation and pricing: Constant regret via Bellman inequalities. \emph{Operations Research}, 69(3), 821--840. \url{https://doi.org/10.1287/opre.2020.2061}
  \bibitem[Wiesemann et al., 2013]{wiesemann2013} Wiesemann, W., Kuhn, D., \& Rustem, B. (2013). Robust Markov decision processes. \emph{Mathematics of Operations Research}, 38(1), 153--183. \url{https://doi.org/10.1287/moor.1120.0566}
  \bibitem[Wu et al., 2024]{wu2024} Wu, Z., Zhang, Y., Li, L., Zhang, Z., Zhao, B., Zhang, Y., \& He, X. (2024). Research on inbound jobs' scheduling in four-way-shuttle-based storage system. \emph{Processes}, 12(1), 223. \url{https://doi.org/10.3390/pr12010223}
  \bibitem[\v{Z}ulj et al., 2022]{zulj2022} \v{Z}ulj, I., Salewski, H., Goeke, D., \& Schneider, M. (2022). Order batching and batch sequencing in an AMR-assisted picker-to-parts system. \emph{European Journal of Operational Research}, 298(1), 182--201. \url{https://doi.org/10.1016/j.ejor.2021.05.033}
\end{thebibliography}

\end{document}
